%% Chapter 3: Methodology: chapters/chapter03.tex

% ==============================================================================
% Chapter Guidelines & Content Description:
% - Theoretical framework and mathematical equations.
% - Computational methods and self-consistent flow (using TikZ).
% - Demonstration of Algorithm floats and Code Listings.
% ==============================================================================

\chapter{Methodology and Theoretical Framework}
\label{ch:methodology}

\section{Density Functional Theory}
The foundation of first-principles calculations is the Hohenberg-Kohn theorems, which state that all ground-state properties of a system are unique functionals of the electron density, $\rho(\mathbf{r})$. The Kohn-Sham single-particle equations are given by:
\begin{equation}
    \left[ -\frac{\hbar^2}{2m_e}\nabla^2 + V_{\text{eff}}(\mathbf{r}) \right] \psi_i(\mathbf{r}) = \epsilon_i \psi_i(\mathbf{r})
    \label{eq:kohn_sham}
\end{equation}
where $V_{\text{eff}}(\mathbf{r})$ is the effective potential defined as:
\begin{equation}
    V_{\text{eff}}(\mathbf{r}) = V_{\text{ext}}(\mathbf{r}) + V_{\text{H}}(\mathbf{r}) + V_{\text{xc}}(\mathbf{r})
    \label{eq:v_eff}
\end{equation}
Here, $V_{\text{ext}}$ is the external potential from nuclei, $V_{\text{H}}$ is the Hartree potential, and $V_{\text{xc}}$ is the exchange-correlation potential.

\section{Self-Consistent Field Loop}
The Kohn-Sham equations must be solved iteratively. Figure~\ref{fig:dft_flowchart} shows the standard self-consistent field (SCF) calculation workflow.

\begin{figure}[H]
    \centering
    \begin{tikzpicture}[
        node distance=1.8cm,
        block/.style={rectangle, draw, fill=blue!10, text width=6cm, text centered, rounded corners, minimum height=1cm},
        decision/.style={diamond, draw, fill=yellow!10, text width=2.5cm, text centered, inner sep=0pt, minimum height=1cm},
        arrow/.style={thick, -{Latex[scale=1.2]}}
    ]
        % Nodes
        \node (start) [block] {1. Initialize trial electron density $\rho_0(\mathbf{r})$};
        \node (pot) [block, below of=start] {2. Calculate effective potential $V_{\text{eff}}(\mathbf{r})$ via Eq.~\eqref{eq:v_eff}};
        \node (ks) [block, below of=pot] {3. Solve Kohn-Sham equations (Eq.~\eqref{eq:kohn_sham}) for $\psi_i$};
        \node (dens) [block, below of=ks] {4. Compute new electron density $\rho_{\text{new}}(\mathbf{r}) = \sum |\psi_i|^2$};
        \node (dec) [decision, below of=dens, yshift=-0.5cm] {Is $|\rho_{\text{new}} - \rho_{\text{old}}| < \eta$?};
        \node (stop) [block, below of=dec, yshift=-0.5cm] {5. Calculation Converged! Output Energy and forces};
        
        % Arrows
        \draw [arrow] (start) -- (pot);
        \draw [arrow] (pot) -- (ks);
        \draw [arrow] (ks) -- (dens);
        \draw [arrow] (dens) -- (dec);
        \draw [arrow] (dec) -- node[anchor=east] {Yes} (stop);
        \draw [arrow] (dec.east) -- ++(2cm,0) |- node[anchor=south, xshift=1.5cm] {No: Mix density} (pot.east);
    \end{tikzpicture}
    \caption{Flowchart of the self-consistent field (SCF) loop in DFT calculations.}
    \label{fig:dft_flowchart}
\end{figure}

\section{Numerical Algorithms}
To optimize structural configurations, we employ a conjugate gradient algorithm. The pseudo-code for our implementation is outlined in Algorithm~\ref{alg:cg}.

\begin{algorithm}[H]
\DontPrintSemicolon
\KwIn{Initial position $\mathbf{x}_0$, tolerance $\epsilon$, max iterations $N$}
\KwOut{Optimized position $\mathbf{x}^*$}
Initialize gradient $\mathbf{g}_0 = \nabla E(\mathbf{x}_0)$ and direction $\mathbf{d}_0 = -\mathbf{g}_0$\;
$k \leftarrow 0$\;
\While{$k < N$ \textbf{and} $\|\mathbf{g}_k\| > \epsilon$}{
  Compute step length $\alpha_k$ to minimize $E(\mathbf{x}_k + \alpha_k \mathbf{d}_k)$\;
  Update position $\mathbf{x}_{k+1} \leftarrow \mathbf{x}_k + \alpha_k \mathbf{d}_k$\;
  Calculate new gradient $\mathbf{g}_{k+1} \leftarrow \nabla E(\mathbf{x}_{k+1})$\;
  Compute scaling factor $\beta_k \leftarrow \frac{\mathbf{g}_{k+1}^T \mathbf{g}_{k+1}}{\mathbf{g}_k^T \mathbf{g}_k}$\;
  Update direction $\mathbf{d}_{k+1} \leftarrow -\mathbf{g}_{k+1} + \beta_k \mathbf{d}_k$\;
  $k \leftarrow k + 1$\;
}
\caption{Conjugate Gradient minimization for structural relaxation.}
\label{alg:cg}
\end{algorithm}

\section{Computational Automation Script}
To automate strain calculations, we developed a Python automation script. The source code is shown in Listing~\ref{lst:python_strain}.

\begin{lstlisting}[style=pythonstyle, caption={Python script to apply biaxial strain to a unit cell file.}, label={lst:python_strain}]
import numpy as np

def apply_biaxial_strain(lattice_vectors, strain_percentage):
    """
    Applies biaxial strain along the a and b directions.
    """
    factor = 1.0 + (strain_percentage / 100.0)
    strained_lattice = np.copy(lattice_vectors)
    strained_lattice[0] *= factor # Strain along a
    strained_lattice[1] *= factor # Strain along b
    return strained_lattice

# Example usage for a 2% tensile strain
lattice = np.array([[3.18, 0.0, 0.0],
                    [-1.59, 2.75, 0.0],
                    [0.0, 0.0, 15.0]])
new_lattice = apply_biaxial_strain(lattice, 2.0)
print("Strained Lattice Vectors:\n", new_lattice)
\end{lstlisting}
